\par\Needspace{8\baselineskip}\bigskip
\pdfbookmark[3]{Exposition, 6}{pallas-II-article6}
\begin{center}6.\end{center}
\noindent\hspace*{\parindent}Nous désignerons l’arc \(XN\) compté de \(X\) vers \(Y\), par \(\AscendingNode\), et la somme de cet arc avec l’arc \(NL\) compté de \(N\) dans le sens du mouvement par \(v\); cette somme est ce que les astronomes nomment la longitude dans l’orbite. Soit \(\tilde\omega\) ce que devient \(v\) lorsque \(L\) devient \(P\), ou la longitude du périhélie; \(v-\tilde\omega\) sera donc l’arc \(PL\) ou l’anomalie vraie. L’inclinaison de l’orbite sur la branche du grand cercle \(XY\) qui va de \(N\) dans le sens direct sera égale à \(ZM\); nous la désignerons par \(i\). Cela posé les 15 formules précédentes \SourceLeaf{454}peuvent être représentées de la manière suivante:
\[
\begin{array}{rl}
42)\quad&\begin{aligned}[t]x&=r(\cos\AscendingNode\mathbin{.}\cos(v-\AscendingNode)\\&\qquad-\sin\AscendingNode\mathbin{.}\sin(v-\AscendingNode)\mathbin{.}\cos i)\end{aligned}\\[5pt]
43)\quad&\begin{aligned}[t]y&=r(\sin\AscendingNode\mathbin{.}\cos(v-\AscendingNode)\\&\qquad+\cos\AscendingNode\mathbin{.}\sin(v-\AscendingNode)\mathbin{.}\cos i)\end{aligned}\\[5pt]
44)\quad&z=r\sin(v-\AscendingNode)\mathbin{.}\sin i
\end{array}
\]
\[
\begin{array}{rl}
45)\quad&A=\BareRoot{mp}\mathbin{.}\sin\AscendingNode\mathbin{.}\sin i\\[4pt]
46)\quad&B=-\BareRoot{mp}\mathbin{.}\cos\AscendingNode\mathbin{.}\sin i\\[4pt]
47)\quad&C=\BareRoot{mp}\mathbin{.}\cos i
\end{array}
\]
\[
\begin{array}{rl}
48)\quad&\begin{aligned}[t]Bz-Cy&=r\BareRoot{mp}\mathbin{.}(-\cos\AscendingNode\mathbin{.}\sin(v-\AscendingNode)\\&\qquad-\sin\AscendingNode\mathbin{.}\cos(v-\AscendingNode)\mathbin{.}\cos i)\end{aligned}\\[5pt]
49)\quad&\begin{aligned}[t]Cx-Az&=r\BareRoot{mp}\mathbin{.}(-\sin\AscendingNode\mathbin{.}\sin(v-\AscendingNode)\\&\qquad+\cos\AscendingNode\mathbin{.}\cos(v-\AscendingNode)\mathbin{.}\cos i)\end{aligned}\\[5pt]
50)\quad&Ay-Bx=r\BareRoot{mp}\mathbin{.}\cos(v-\AscendingNode)\mathbin{.}\sin i
\end{array}
\]
\[
\begin{array}{rl}
51)\quad&\begin{aligned}[t]CG-BH&=mpe(\cos\AscendingNode\mathbin{.}\cos(\tilde\omega-\AscendingNode)\\&\qquad-\sin\AscendingNode\mathbin{.}\sin(\tilde\omega-\AscendingNode)\mathbin{.}\cos i)\end{aligned}\\[5pt]
52)\quad&\begin{aligned}[t]AH-CF&=mpe(\sin\AscendingNode\mathbin{.}\cos(\tilde\omega-\AscendingNode)\\&\qquad+\cos\AscendingNode\mathbin{.}\sin(\tilde\omega-\AscendingNode)\mathbin{.}\cos i)\end{aligned}\\[5pt]
53)\quad&BF-AG=mpe\sin(\tilde\omega-\AscendingNode)\mathbin{.}\sin i
\end{array}
\]
\[
\begin{array}{rl}
54)\quad&\begin{aligned}[t]F&=e\BareRoot{mp}\mathbin{.}(-\cos\AscendingNode\mathbin{.}\sin(\tilde\omega-\AscendingNode)\\&\qquad-\sin\AscendingNode\mathbin{.}\cos(\tilde\omega-\AscendingNode)\mathbin{.}\cos i)\end{aligned}\\[5pt]
55)\quad&\begin{aligned}[t]G&=e\BareRoot{mp}\mathbin{.}(-\sin\AscendingNode\mathbin{.}\sin(\tilde\omega-\AscendingNode)\\&\qquad+\cos\AscendingNode\mathbin{.}\cos(\tilde\omega-\AscendingNode)\mathbin{.}\cos i)\end{aligned}\\[5pt]
56)\quad&H=e\BareRoot{mp}\mathbin{.}\cos(\tilde\omega-\AscendingNode)\mathbin{.}\sin i.
\end{array}
\]
\par Multipliant les équations 11, 12, 13 par \(A,B,C\) et ajoutant, on trouve
\[
x(BH-CG)+y(CF-AH)+z(AG-BF)=mp(r-p).
\]
\par Or le premier membre de cette équation dévient, moyennant les équations 27—29, 36—38
\[
=-mper\cos PL=-mper\cos(v-\tilde\omega);
\]
ainsi on a
\[
\begin{array}{rl}57)\quad&r=\dfrac{p}{1+e\cos(v-\tilde\omega)}\mathpunct{\cdot}\end{array}
\]
\par De même l’équation 20 donne moyennant 27—29, 39—41
\[
er\BareRoot{mp}\mathbin{.}\cos\Pi L=\dfrac{pr\mathrm{d}r}{\mathrm{d}t}
\]
ou à cause de
\[
\Pi L=v-\tilde\omega-90^{0}:
\]
\[
\begin{array}{rl}58)\quad&\dfrac{\mathrm{d}r}{\mathrm{d}t}=e\sin(v-\tilde\omega)\mathbin{.}\BareRoot{\dfrac{m}{p}}\mathpunct{\cdot}\end{array}
\]
\par La différentiation de l’équation 57 donne \(\dfrac{p\mathrm{d}r}{rr\mathrm{d}t}=e\sin(v-\tilde\omega)\mathbin{.}\dfrac{\mathrm{d}v}{\mathrm{d}t}\); \SourceLeaf{455}combinant ce résultat avec l’équation 58 on a
\[
\begin{array}{rl}59)\quad&\dfrac{rr\mathrm{d}v}{\mathrm{d}t}=\BareRoot{mp}.\end{array}
\]
\par On peut remarquer que \(\tfrac12rr\mathrm{d}v\) est l’élément de l’aire décrite autour du soleil dans l’élément du tems; or on obtient les projections de cette aire aux trois plans fondamentaux en la multipliant par \(\cos XM\), \(\cos YM\), \(\cos ZM\): donc \(\tfrac12A\mathrm{d}t\), \(\tfrac12B\mathrm{d}t\), \(\tfrac12C\mathrm{d}t\) doivent être égaux à ces mêmes projections, ce qui au reste se vérifie de soi-même.
\par Remarquons enfin qu’on a
\[
\begin{aligned}
r\sin(v-\tilde\omega)&=\BareRoot{\dfrac{p}{m}}\mathbin{.}\dfrac{r\mathrm{d}r}{e\mathrm{d}t}\\
&=\BareRoot{ap}\mathbin{.}\sin E=a\BareRoot{(1-ee)}\mathbin{.}\sin E
\end{aligned}
\]
\[
\begin{aligned}
r+r\cos(v-\tilde\omega)&=\dfrac{p-r(1-e)}{e}\\
&=\dfrac{p-a(1-e)(1-e\cos E)}{e}=a(1-e)(1+\cos E)
\end{aligned}
\]
\[
\begin{aligned}
r-r\cos(v-\tilde\omega)&=\dfrac{-p+r(1+e)}{e}\\
&=\dfrac{-p+a(1+e)(1-e\cos E)}{e}=a(1+e)(1-\cos E).
\end{aligned}
\]
\par Divisant la première de ces équations par la seconde, ou la troisième par la première, il vient
\[
\begin{array}{rl}60)\quad&\operatorname{tang}\tfrac12(v-\tilde\omega)=\BareRoot{\dfrac{1+e}{1-e}}\mathbin{.}\operatorname{tang}\tfrac12E.\end{array}
\]
\par\medskip\begin{center}\rule{.18\textwidth}{.3pt}\end{center}
